\documentclass[11pt]{article}
%\textwidth17cm \textheight23cm
%\hoffset-15mm \voffset-19mm
\parskip4pt plus1pt minus1pt

\usepackage[T1]{fontenc}
\usepackage[utf8]{inputenc}
\usepackage{lmodern}
\usepackage{microtype}
\usepackage[margin=1in]{geometry}
\usepackage{amsmath,amssymb,amsthm,mathtools}
\usepackage{booktabs}
\usepackage{enumitem}
\usepackage{xcolor}
\usepackage{hyperref}
\usepackage{mathrsfs}

\hypersetup{colorlinks=true,linkcolor=blue,citecolor=blue,urlcolor=blue}

\newtheorem{theorem}{Theorem}[section]
\newtheorem{lemma}[theorem]{Lemma}
\newtheorem{proposition}[theorem]{Proposition}
\newtheorem{corollary}[theorem]{Corollary}
\newtheorem{conjecture}[theorem]{Conjecture}
\newtheorem{definition}[theorem]{Definition}
\newtheorem{remark}[theorem]{Remark}
\newtheorem{example}[theorem]{Example}




\newcommand{\R}{\mathbb{R}}
\newcommand{\tr}{\operatorname{tr}}
\newcommand{\Spec}{\operatorname{Spec}}
\newcommand{\rank}{\operatorname{rank}}
\newcommand{\im}{\operatorname{im}}
\newcommand{\Ker}{\operatorname{ker}}
\newcommand{\Span}{\operatorname{span}}
\newcommand{\supp}{\operatorname{supp}}
\newcommand{\bd}{\partial}
\newcommand{\ip}[2]{\left\langle #1,#2\right\rangle}
\newcommand{\norm}[1]{\left\lVert #1\right\rVert}
\newcommand{\diag}{\operatorname{diag}}
\newcommand{\one}{\mathbf 1}

\newcommand{\calC}{\mathcal{C}}
\newcommand{\calR}{\mathcal{R}}
\newcommand{\calF}{\mathcal{F}}
\newcommand{\calK}{\mathcal{K}}
\newcommand{\calN}{\mathcal{N}}
\newcommand{\scrH}{\mathscr{H}}
\newcommand{\scrU}{\mathscr{U}}
\newcommand{\scrV}{\mathscr{V}}
\newcommand{\scrE}{\mathscr{E}}

\newcommand{\blue}[1]{{\color{blue}#1}}

%==================================
\begin{document}

\title{\bf The Duval--Reiner Conjecture:  Counterexamples and the  Second Partial-Sum Inequality}

\author{
Jing Huang\footnote{Email address:
jhuangmath@gzhu.edu.cn
}}

\date{\small
School of Mathematics and Information Science, Guangzhou University, Guangzhou,  510006, China\\
}


\maketitle

\begin{abstract}
Let \(F\subseteq\binom{V}{q}\) be a \(q\)-uniform family on a finite
vertex set \(V\). Write \(s_r(F)\) for the sum of the \(r\) largest
eigenvalues of its simplicial up-Laplacian and \(d_F(v)\) for the
degree of \(v\in V\). Then $D_r(F)=\sum_{v\in V}\min\{d_F(v),r\}$
is the \(r\)-th partial sum of the conjugate degree sequence of \(F\).
The majorization assertion in the Duval--Reiner conjecture
[Trans. Amer. Math. Soc., 2002]
states that \(s_r(F)\le D_r(F)\) for every \(q\)-uniform family \(F\)
and every \(r\ge1\).
We disprove this assertion in two complementary senses: for every
\(r\ge5\), there is a strict counterexample at index \(r\) in some
uniformity, while every uniformity \(q\ge3\) admits a strict
counterexample at some index \(r\ge5\).
In contrast, we prove the universal inequality
\(s_2(F)\le D_2(F)\) and classify all equality cases.
The counterexamples are obtained from two \(3\)-uniform seeds with
explicitly computed characteristic polynomials through
defect-preserving ridge-whiskering and set-complement duality.
For the second partial sum, core completion reduces the problem to the
boundary matrix of a complete simplex, where Ky Fan variational and
compression arguments yield both the inequality and its equality
classification.

\medskip
{\bf Key words}: Duval--Reiner conjecture, simplicial Laplacian, spectral majorization,
conjugate degree sequence, equality classification.

\medskip
{\bf MSC}:  05E45, 05C65, 05C50, 15A42.

\end{abstract}

\section{Introduction}
\label{sec:introduction}
Let \(V\) be a finite set and let \(q\geq1\) be an integer. We write
\(\binom{V}{q}\) for the collection of all \(q\)-element subsets of \(V\).
A subfamily \(F\subseteq\binom{V}{q}\) is called a
\emph{\(q\)-uniform family}; its members are called \emph{facets}. Its
downward closure $\langle F\rangle=\{\tau\subseteq V:\tau\subseteq\sigma
\text{ for some }\sigma\in F\}$
is the \emph{simplicial complex} generated by \(F\).  For a face \(\tau\in\langle F\rangle\), define $\dim\tau=|\tau|-1.$
A \emph{ridge} is a face of dimension \(q-2\), equivalently a
\((q-1)\)-element subset of a facet.
Fix a linear order on \(V\).  For
\(\sigma=\{v_0<\cdots<v_{q-1}\}\), define the signed boundary by
\[
	\bd e_\sigma
	=
	\sum_{j=0}^{q-1}(-1)^j
	e_{\sigma\setminus\{v_j\}}.
\]
Let \(B_F\) be the corresponding boundary matrix, with columns indexed by
the facets of \(F\) and rows indexed by the ridges occurring in those
facets.  For \(q=1\), the unique possible row is indexed by
\(\varnothing\).  Let
$
L_F^+=B_FB_F^{\mathsf T}
$
and
$
M_F=B_F^{\mathsf T}B_F.
$
Thus \(L_F^+\) is the simplicial up-Laplacian, whereas \(M_F\) is its
facet Gram matrix. Both matrices are real symmetric positive semidefinite, and they have the same positive eigenvalues, with multiplicity.
List them in weakly decreasing order and pad the list with zeros when
necessary:
\[
	\lambda_1(F)\geq\lambda_2(F)\geq\cdots\geq0,
	\qquad
	s_r(F)=\sum_{i=1}^r\lambda_i(F).
\]
Thus \(s_r(F)\) is the \(r\)-th Ky Fan sum of the simplicial
up-Laplacian.
For \(v\in V\), let $d_F(v)=|\{\sigma\in F:v\in\sigma\}|$
be its vertex degree.  The conjugate vertex-degree sequence and its \(r\)-th partial sum are
\[
	d_i^T(F)=|\{v\in V:d_F(v)\geq i\}|,
	\qquad
	D_r(F)=\sum_{i=1}^r d_i^T(F)
	=
	\sum_{v\in V}\min\{d_F(v),r\}.
\]
Both sequences have total sum \(q|F|\):
\[
	\sum_i\lambda_i(F)
	=
	\tr(L_F^+)
	=
	q|F|
	=
	\sum_{v\in V}d_F(v)
	=
	\sum_i d_i^T(F).
\]

This framework extends the graph setting. When \(q=2\), \(F\) is the
edge set of a finite simple graph, \(B_F\) is an oriented vertex--edge
incidence matrix, and \(L_F^+=B_FB_F^{\mathsf T}\) is its usual
Laplacian with isolated vertices omitted. The graph Laplacian goes back
to Kirchhoff's work on electrical networks \cite{Kirchhoff1847}, while
Eckmann developed the corresponding finite-dimensional Hodge theory for
simplicial complexes \cite{Eckmann1944}.
Beyond these origins, their spectral properties are of considerable combinatorial interest.
For graphs, Grone and Merris conjectured
that the graph Laplacian spectrum is majorized by the conjugate degree
sequence \cite{GroneMerris1994}; Duval and Reiner subsequently formulated a higher-dimensional generalization
for
uniform simplicial families
\cite[Conjecture~1.2]{DuvalReiner2002}.
To state their conjecture in full, recall that a \(q\)-family on
\(V=\{v_1<\cdots<v_n\}\) is \emph{shifted} if, whenever
\(1\le i_1<\cdots<i_q\le n\) and
\(1\le j_1<\cdots<j_q\le n\) satisfy
\(i_\ell\le j_\ell\) for every \(1\le\ell\le q\), the inclusion
$
\{v_{j_1},\ldots,v_{j_q}\}\in F
$
implies
$
\{v_{i_1},\ldots,v_{i_q}\}\in F.
$
A family is \emph{isomorphic to a shifted family} if some relabeling of
its vertices makes it shifted.
\begin{conjecture}
\label{conj:duval-reiner}
For every \(q\)-uniform family \(F\) and every \(r\ge1\), $s_r(F)\leq D_r(F).$
Moreover, equality holds simultaneously for all \(r\ge1\) if and only if
\(F\) is isomorphic to a shifted family.
\end{conjecture}

Conceptually, the conjecture asks whether vertex-incidence data control
every leading Ky Fan sum of a higher-dimensional boundary operator, with
shifted families as the exact model.
Some recent papers use the name ``Duval--Reiner conjecture'' for the
majorization assertion alone \cite{HanLuWang2026,Lew2025}. In the
original formulation, however, the equality clause concerns coincidence
of the entire spectral and conjugate vertex-degree sequences, rather
than equality of a single partial sum \(s_r(F)=D_r(F)\).
Beyond these nuances, this framework has influenced proposed generalizations of Brouwer's
spectral-sum conjecture to simplicial complexes \cite{Abebe12,Abebe}.


Duval and Reiner proved the first partial-sum inequality in every
uniformity and the second in uniformity \(q=2\); they also established
invariance under set and family complementation and preservation under
joins \cite{DuvalReiner2002}. Duval later extended the exact
shifted spectral identity to pairs of shifted complexes
\cite{Duval2005}. In the graph case, related results treated trees
\cite{Stephen2007}, one-regular semibipartite graphs \cite{Katz2005},
cographs \cite{BapatLalPati2008}, and near-threshold graphs
\cite{Kirkland2009}. Stephen also reduced the relative graph-pair
majorization problem to the ordinary graph case \cite{Stephen2007}, and
Bai's proof of the Grone--Merris conjecture completed the majorization
theorem for all graphs \cite{Bai2011}.

For higher uniformities, Fan, Wu, and Wang obtained local degree bounds
and balancedness criteria at the first index
\cite{FanWuWang2025}. Lew proved a general Ky Fan bound in terms of ridge
degrees and verified all the Duval--Reiner inequalities for
\(q\)-partite \(q\)-families \cite{Lew2025}. Zhang and Fan refined the
universal first-index bound and established homological and local
equality criteria \cite{ZhangFan2026}. Most recently, Han, Lu, and Wang
proved the second partial-sum inequality for arbitrary \(3\)-families
\cite{HanLuWang2026}. Thus, before the present work, the second
inequality remained open for unrestricted \(q\)-families with \(q\ge4\),
and no complete majorization result was known in any unrestricted
uniformity \(q\ge3\).


Our first main result gives strict counterexamples in two complementary
quantifier regimes, thereby disproving the majorization assertion and
hence the Duval--Reiner conjecture.
\begin{theorem}
\label{thm:counterexamples}
For every integer \(r\ge5\), there exist an integer \(q\ge3\), a finite
set \(V\), and a family \(F\subseteq\binom{V}{q}\) such that
$s_r(F)> D_r(F).$
Moreover, every integer \(q\ge3\) admits such a counterexample at some
index \(r\ge5\). In both assertions, \(F\) may be chosen so that
$\bigcap_{\sigma\in F}\sigma=\varnothing.$
\end{theorem}
Exact characteristic polynomials certify two
\(3\)-uniform counterexample seeds, whose positive defects are propagated
by ridge-whiskering and set-complement duality.
The order of the quantifiers is essential: the theorem does not assert
the existence of a counterexample for every pair \((q,r)\) with
\(q\ge3\) and \(r\ge5\). The empty-intersection condition rules out
counterexamples obtained merely by coning a lower-uniformity example.
Our second main result establishes the second partial-sum inequality in
every uniformity and gives a complete classification of its equality
cases.
To formulate the equality conditions, let
\[
P(F)=\{v\in V:d_F(v)=1\},
\qquad
Q(F)=\{v\in V:d_F(v)\ge2\}
\]
be the pendant and core vertex sets. For \(q\ge2\), put
\[
Q=Q(F),\qquad
m=|Q|,\qquad
\calK=F\cap\binom Qq,\qquad
\calN=\binom Qq\setminus\calK,
\]
and
\[
\calC=
\left\{
\rho\in\binom Q{q-1}:
\rho\cup\{z\}\in F
\text{ for some }z\in P(F)
\right\}.
\]
Let \(\bd_Q\) be the signed boundary matrix of \(\binom Qq\), and, for
\(\mathcal A\subseteq\binom Qq\), let \(\bd_{\mathcal A}\) be the column
submatrix of \(\bd_Q\) indexed by \(\mathcal A\).

\begin{theorem}
	\label{thm:main-inequality}
	For every integer \(q\ge1\), every finite set \(V\), and every family
\(F\subseteq\binom{V}{q}\), we have $s_2(F)\le D_2(F).$
	After isolated vertices are deleted, equality holds if and only if at
	least one of the following conditions is satisfied.
	\begin{enumerate}[label=\textup{(\roman*)},leftmargin=*]
		\item \(q=1\), with \(F\) arbitrary.

		\item \(q\geq2\) and \(F\) has at most two facets.

		\item \(q\geq2\), \(m=q-1\), and \( F=\{Q\cup\{z\}:z\in Z\} \)  for some \(Z\subseteq V\setminus Q\) with \(|Z|\geq2\).

		\item \(q\ge2\), every facet of \(F\) contains at most one pendant
		vertex, and all of the following hold:
		\begin{enumerate}[label=\textup{(\alph*)},leftmargin=2.4em]
			\item \(m\ge q+1\), \(|\calC|\le2\), and
			\(
				\rank_{\R}\bd_{\calN}
				\le\binom{m-1}{q-1}-2;
			\)
			\item \(|\rho\cap\gamma|\le q-3\) for all distinct
			\(\rho,\gamma\in\calC\);
			\item \(\rho\cup\{x\}\in\calK\) for every
			\(\rho\in\calC\) and \(x\in Q\setminus\rho\).
		\end{enumerate}
	\end{enumerate}
\end{theorem}
This theorem classifies equality only at \(r=2\), not equality of the
full sequences in Conjecture~\ref{conj:duval-reiner}; notably,
case~\textup{(iv)} allows \(\calN\neq\varnothing\), so equality does not
require a complete core. For \(r=2\), core
completion and Ky Fan's variational principle reduce the spectral problem
to the complete-simplex boundary block; equality is governed jointly by
the rank of the missing-core boundary matrix and the configuration of the
pendant-supporting ridges.

Since the inequalities at \(r=1,2\) hold in every uniformity, whereas
Theorem~\ref{thm:counterexamples} provides a counterexample at every
index \(r\ge5\), the only unresolved universal majorization inequalities
are $s_3(F)\le D_3(F)$ and $s_4(F)\le D_4(F)$
for every finite uniform family \(F\).

The paper is organized as follows.
Section~\ref{sec:notation} develops the spectral and boundary-map
notation and the required matrix tools.
Section~\ref{sec:seeds} proves Theorem~\ref{thm:counterexamples} from two
\(3\)-uniform seeds using ridge-whiskering and set-complement duality.
Section~\ref{sec:core-inequality} proves
Theorem~\ref{thm:main-inequality} and classifies all equality cases.
The appendix records the seed matrices and verifies their characteristic
polynomials by exact arithmetic.
	
	\section{Preliminary}\label{sec:notation}

We retain the notation from the Introduction. For every nonempty
oriented face \(\sigma=\{v_0<\cdots<v_k\}\), extend the boundary
convention by
\[
\bd e_\sigma
=
\sum_{i=0}^k(-1)^i e_{\sigma\setminus\{v_i\}},
\qquad
\bd e_\varnothing=0.
\]
For faces \(\rho\) and \(\sigma\), define the oriented incidence
coefficient by
\[
\varepsilon(\rho,\sigma)
=
\ip{e_\rho}{\bd e_\sigma}
=
\begin{cases}
(-1)^i,
&\rho=\sigma\setminus\{v_i\}\text{ for some }i,\\
0,
&\text{otherwise}.
\end{cases}
\]
The identity \(\bd^2=0\) follows from the usual cancellation: for
\(0\le a<b\le k\), the two contributions to the coefficient of
\(e_{\sigma\setminus\{v_a,v_b\}}\) in \(\bd^2e_\sigma\) are
$(-1)^a(-1)^{b-1}$ and $(-1)^b(-1)^a$
whose sum is zero.
When convenient, we index the rows of \(B_F\) by all of
\(\binom{V}{q-1}\), adjoining zero rows for the ridges not occurring in
\(F\). Changing the vertex order or the oriented face bases replaces
\(B_F\) by \(PB_FQ\) for signed permutation matrices \(P\) and \(Q\).
Consequently, \(L_F^+\) and \(M_F\) change only by orthogonal similarity,
so their spectra are independent of these choices.

For a real symmetric matrix \(A\), write \(A\succeq0\) when $
\mathbf x^{\mathsf T}A\mathbf x\ge0$  for every $\mathbf x,$
and write \(A\preceq0\) when \(-A\succeq0\). If
\(A\in\R^{n\times n}\), let \(\Spec(A)\) denote its eigenvalue multiset,
and arrange its eigenvalues as
$\lambda_1(A)\ge\cdots\ge\lambda_n(A).$
For \(0\le r\le n\), define
\[
s_r(A)=\sum_{i=1}^r\lambda_i(A),
\qquad
s_0(A)=0.
\]
For a positive-semidefinite matrix, append zero eigenvalues and thereby
extend \(s_r(A)\) to every \(r\ge0\). With \(s_0(F)=0\), the notation
from the Introduction satisfies $s_r(F)=s_r(L_F^+)=s_r(M_F)\  (r\ge0).$
For \(a\in\R\), put \(a_+=\max\{a,0\}\). For a coordinate vector
\(\mathbf x=(x_\omega)_{\omega\in\Omega}\) and a subspace
\(\scrU\subseteq\R^\Omega,\)  define
\[
\supp(\mathbf x)=\{\omega\in\Omega:x_\omega\ne0\},
\qquad
\supp(\scrU)=\bigcup_{\mathbf x\in\scrU}\supp(\mathbf x).
\]

The case \(q=1\) is immediate. If \(F=\varnothing\), then
\(s_r(F)=D_r(F)=0\). If \(F\ne\varnothing\), then \(B_F\) is the
one-row all-one matrix, so \(L_F^+\) has the single positive eigenvalue
\(|F|\), while every nonisolated vertex has degree one. Hence
$s_r(F)=D_r(F)=|F|$ for \(r\ge1\).
Deleting isolated vertices changes neither side. Accordingly, except
where stated otherwise, we henceforth assume
$q\ge2$ and $V=\bigcup_{\sigma\in F}\sigma.$
	
We next collect the variational and compression tools used in the proof
of Theorem~\ref{thm:main-inequality}.   If \(\scrU\) is a subspace
	of a finite-dimensional Euclidean space, let \(\Pi_{\scrU}\) denote the
	orthogonal projection onto \(\scrU\).  We use \(\norm{\cdot}\) for the
	Euclidean norm on vectors and its induced spectral norm on matrices.  For a
	finite set \(\Omega\), a rank-\(k\)
	orthogonal projection on \(\R^\Omega\) is a matrix \(\Pi\in\R^{\Omega\times\Omega}\)
	such that \(\Pi^2=\Pi=\Pi^{\mathsf T}\) and \(\rank(\Pi)=k\).
The following is Ky Fan's maximum principle; while standard,
its equality case requires explicit formalization for our application.
	\begin{lemma}\label{lem:kyfan}
		Let \(A\in\R^{\Omega\times\Omega}\) be symmetric.
		Then, for every \(1\le k\le |\Omega|\),
		\[
		s_k(A)
		=
		\max_\Pi \tr(\Pi A),
		\]
		where \(\Pi\) ranges over all rank-\(k\) orthogonal projections on \(\R^\Omega\).
		Moreover, if \(a=\lambda_k(A)\) is the \(k\)-th largest eigenvalue of \(A\), a \(k\)-dimensional subspace \(\scrU\)  satisfies \(\tr(\Pi_{\scrU}A)=s_k(A)\) if and only if $\scrE_{>a}\subseteq\scrU\subseteq\scrE_{>a}\oplus\scrE_{=a},$
		where \(\scrE_{>a}\) and \(\scrE_{=a}\) denote the sums of the eigenspaces
		of \(A\) with eigenvalues strictly greater than, and equal to \(a\),
		respectively.
	\end{lemma}
	
	\begin{proof}
		The variational formula is Ky Fan's maximum principle \cite[Theorem~1]{Fan}.
		It remains to describe the equality case.
		Put \(n=|\Omega|\), and let \(\lambda_1\ge\cdots\ge\lambda_n\) be the eigenvalues of \(A\).
		Choose an orthonormal eigenbasis \(\mathbf u_1,\ldots,\mathbf u_n\)
		with \(A\mathbf u_j=\lambda_j\mathbf u_j\). For a \(k\)-dimensional subspace
		\(\scrU\), put \(a_j=\norm{\Pi_{\scrU}\mathbf u_j}^2\). Then
		\[
		s_k(A)=\sum_{\lambda_j>a}\lambda_j+ \bigl(k-\dim\scrE_{>a}\bigr)a.
		\]
		Because \(\sum_j a_j=k\), subtraction gives the exact identity
		\[
		s_k(A)-\tr(\Pi_{\scrU}A)
		=\sum_{\lambda_j>a}(\lambda_j-a)(1-a_j)
		+\sum_{\lambda_j<a}(a-\lambda_j)a_j.
		\]
		Every summand is nonnegative. Equality therefore holds exactly when
		\(a_j=1\) for \(\lambda_j>a\) and \(a_j=0\) for \(\lambda_j<a\).
		For an orthogonal projection, \(\norm{\Pi_{\scrU}\mathbf u_j}=1\) is equivalent to
		\(\mathbf u_j\in\scrU\), and \(\norm{\Pi_{\scrU}\mathbf u_j}=0\) is equivalent to
		\(\mathbf u_j\perp\scrU\). This is precisely the displayed containment.
	\end{proof}
We will need the following standard monotonicity property of Ky Fan sums.
	
	\begin{lemma}[\cite{Bhatia}]\label{lem:kyfan-monotone}
		Let \(A,B\in\R^{n\times n}\) be symmetric. If \(B-A\succeq 0\), then
		\(s_r(A) \le s_r(B)\)
		for every \(1\le r\le n\).
	\end{lemma}
	
	For a real symmetric matrix
	\(A=\sum_j\lambda_j\mathbf u_j\mathbf u_j^{\mathsf T}\), define its
	\emph{positive part} by
	$A_+=\sum_{\lambda_j>0}\lambda_j\mathbf u_j\mathbf u_j^{\mathsf T}.$
	Thus \(\tr(A_+)=\sum_j(\lambda_j(A))_+\), independently of the
	chosen orthonormal eigenbasis.
We also record a subadditivity property for the trace of the positive part.
	
	\begin{lemma}\label{lem:positive-part}
		Let \(A_1,\ldots,A_\ell\in\R^{\Omega\times\Omega}\) be symmetric,
		where \(\Omega\) is finite. Then
		\[
		\tr\bigl((A_1+\cdots+A_\ell)_+\bigr)
		\le
		\sum_{i=1}^\ell \tr((A_i)_+).
		\]
	\end{lemma}
	
	\begin{proof}
		For an arbitrary real symmetric matrix
		\(A=\sum_j\lambda_j(A)\mathbf u_j\mathbf u_j^{\mathsf T}\) and any orthogonal
		projection \(\Pi\), the spectral theorem gives
		\[
		\tr(\Pi A)=\sum_j\lambda_j(A)\norm{\Pi\mathbf u_j}^2
		\le\sum_{\lambda_j(A)>0}\lambda_j(A)=\tr(A_+).
		\]
		Equality is attained by the projection onto the direct sum of the
		positive eigenspaces, and by \(\Pi=0\) if \(A\preceq0\).  Therefore
		\(\tr(A_+)=\max_\Pi\tr(\Pi A)\), where \(\Pi\) ranges over projections of
		all ranks.  Applying this identity to the sum gives
		\begin{align*}
			\tr\bigl((A_1+\cdots+A_\ell)_+\bigr)
			&=\max_\Pi\sum_{i=1}^\ell\tr(\Pi A_i)\\
			&\le\sum_{i=1}^\ell\max_{\Pi_i}\tr(\Pi_iA_i)
			=\sum_{i=1}^\ell\tr((A_i)_+).
		\end{align*}
	\end{proof}

To handle principal submatrices, we require the following interlacing bound.
We state the equality condition explicitly in terms of index support to facilitate our structural classification later.
	
	\begin{lemma}\label{lem:compression-equality}
		Let \(B\succeq0\) be a real symmetric matrix, and let \(A\) be a principal submatrix of \(B\).
		For every \(1\le r\le\dim A\), $s_r(A)\le s_r(B).$
		Equality holds if and only if there exists an \(r\)-dimensional subspace \(\scrU\) supported entirely on the indices of \(A\) such that \(\tr(\Pi_{\scrU}B)=s_r(B)\).
	\end{lemma}
	
	\begin{proof}
		Up to a permutation of coordinates, we can write \(B\) in block form as
		\[
		B = \begin{pmatrix} A & C \\ C^{\mathsf T} & D \end{pmatrix}.
		\]
		By Lemma~\ref{lem:kyfan}, \(s_r(A) = \max_U \tr(U^{\mathsf T} A U)\), where \(U\) ranges over all matrices with \(r\) orthonormal columns of the same row dimension as \(A\).
		For any such \(U\), the extended block matrix \(V = \begin{pmatrix} U \\ 0 \end{pmatrix}\) also has \(r\) orthonormal columns. Matrix multiplication yields
$\tr(U^{\mathsf T} A U) = \tr(V^{\mathsf T} B V).$
Applying Lemma~\ref{lem:kyfan} to \(B\) gives \(\tr(V^{\mathsf T} B V) \le s_r(B)\). Taking the maximum over all admissible \(U\) yields \(s_r(A) \le s_r(B)\).
		
		Equality holds exactly when the maximum is attained by some extended matrix \(V\) whose lower block is zero. The column space of such a \(V\) is an \(r\)-dimensional subspace \(\scrU\) supported entirely on \(A\), and it satisfies \(\tr(\Pi_{\scrU}B)=s_r(B)\).
	\end{proof}
	

We now present a rigid trace bound for a specific class of symmetric block matrices, which models our boundary operators.
\begin{lemma}\label{lem:positive-excess-rigidity}
Let \(a, b_1, \ldots, b_\ell > 0\), and let \(\mathbf z_1, \ldots, \mathbf z_\ell\) be vectors of squared norm \(a\) in a finite-dimensional real inner-product space \(\scrH_0\). Define \(X:\R^\ell\to\scrH_0\) on the standard basis by \(Xe_i=\sqrt{b_i}\,\mathbf z_i\), and set
		\[
		T=
		\begin{pmatrix}
			0&X\\
			X^{\mathsf T}&\diag(b_1,\ldots,b_\ell)-aI_\ell
		\end{pmatrix}
		\quad\text{on }\scrH_0\oplus\R^\ell.
		\]
		Then \(\tr(T_+)\le\sum_{i=1}^{\ell}b_i\). Equality holds if and only if the vectors \(\mathbf z_i\) are pairwise orthogonal. In this case, the nonzero eigenvalues of \(T\) are \(b_1,\ldots,b_\ell\) with respective eigenvectors \((\mathbf z_i,\sqrt{b_i}\,e_i)\), and \(-a\) with multiplicity \(\ell\).
	\end{lemma}
	
	\begin{proof}
		For \(1\le i\le\ell\), define \(X_i:\R^\ell\to\scrH_0\) by \(X_i e_j=\delta_{ij}\sqrt{b_i}\,\mathbf z_i\), and put
		\[
		T_i=
		\begin{pmatrix}
			0&X_i\\
			X_i^{\mathsf T}&(b_i-a)e_ie_i^{\mathsf T}
		\end{pmatrix}.
		\]
		Then \(T=\sum_iT_i\). In an orthonormal basis of \(\scrH_0\) extending \(\mathbf z_i/\sqrt{a}\), the principal block of \(T_i\) is orthogonally similar to
		\[
		\begin{pmatrix}
			0&\sqrt{ab_i}\\
			\sqrt{ab_i}&b_i-a
		\end{pmatrix}.
		\]
		The nonzero eigenvalues of \(T_i\) are exactly the roots of \((\lambda-b_i)(\lambda+a)\), namely \(b_i\) and \(-a\), with orthogonal eigenvectors
		\[
		\mathbf u_i=(\mathbf z_i,\sqrt{b_i}\,e_i),
		\qquad
		\mathbf v_i=(-\sqrt{b_i/a}\,\mathbf z_i,\sqrt a\,e_i).
		\]
		Both have squared norm \(a+b_i\). Spectral decomposition gives
		\[
		T_i=\frac{b_i}{a+b_i}\mathbf u_i\mathbf u_i^{\mathsf T}
		-\frac{a}{a+b_i}\mathbf v_i\mathbf v_i^{\mathsf T}.
		\]
		Thus \(\tr((T_i)_+)=b_i\), and Lemma~\ref{lem:positive-part} yields
		\[
		\tr(T_+)\le\sum_{i=1}^{\ell}\tr((T_i)_+)
		=\sum_{i=1}^{\ell}b_i.
		\]
		If equality holds, let \(\Pi\) be the orthogonal projection onto the positive eigenspace of \(T\). Since
		\[
		\tr(\Pi T_i)
		=\frac{b_i}{a+b_i}\norm{\Pi\mathbf u_i}^{2}
		-\frac{a}{a+b_i}\norm{\Pi\mathbf v_i}^{2}\le b_i,
		\]
		equality forces \(\Pi\mathbf u_i=\mathbf u_i\) and \(\Pi\mathbf v_i=0\) for all \(i\). Using \(\Pi=\Pi^{\mathsf T}\), for \(i\ne j\),
		\[
		0=\ip{\mathbf u_i}{\Pi\mathbf v_j}=\ip{\Pi\mathbf u_i}{\mathbf v_j}=\ip{\mathbf u_i}{\mathbf v_j}
		=-\sqrt{b_j/a}\,\ip{\mathbf z_i}{\mathbf z_j}.
		\]
		Hence, the vectors \(\mathbf z_i\) are pairwise orthogonal.
		Conversely, if the \(\mathbf z_i\) are pairwise orthogonal, the subspaces
		\[
			\Span\{(\mathbf z_i,0),(0,e_i)\}
			\qquad(1\le i\le\ell)
		\]
		are mutually orthogonal and \(T\)-invariant.  On the \(i\)-th subspace,
		\(T\) has eigenvalues \(b_i\) and \(-a\), and \(T=0\) on their
		orthogonal complement.  The assertions follow.
	\end{proof}



We conclude this section with a simple algebraic bound for eigenvalue partial sums.
	\begin{lemma}\label{lem:excess-controls}
		Let \(A\succeq0\) and \(a\ge0\). For every \(r\ge1\),
		\[
		s_r(A)\le ra+\sum_{\lambda\in\Spec(A)}(\lambda-a)_+.
		\]
	\end{lemma}
	
	\begin{proof}
		Let \(\lambda_1\ge\lambda_2\ge\cdots\ge0\) be the eigenvalues of \(A\), extended by zeros. Since \(\lambda_j\le a+(\lambda_j-a)_+\) for all \(j\),
		\[
		s_r(A)
		=
		\sum_{j=1}^r\lambda_j
		\le
		ra+\sum_{j=1}^r(\lambda_j-a)_+
		\le
		ra+\sum_{\lambda\in\Spec(A)}(\lambda-a)_+.
		\]
	\end{proof}
	
\section{Systematic Counterexamples}\label{sec:seeds}
This section develops the systematic counterexamples required for Theorem \ref{thm:counterexamples}.
For a uniform family \(F\) and an integer \(r\ge1\), define its signed
defect at index \(r\) by
\[
\Delta_r(F)=s_r(F)-D_r(F).
\]
Thus \(F\) is a strict counterexample at index \(r\) precisely when
\(\Delta_r(F)>0\).
Two $3$-uniform seeds, identified through an exhaustive computer search on seven vertices,
play complementary roles.
  A sixteen-facet seed
	gives the counterexample at \(r=5\) directly.  A fifteen-facet seed drives
	the transfer construction: complementation sends its defect at \(r=9\) to
	\(15-9=6\), while \(t\) ridge whiskers preserve its defect at \(r=8\) and
	complementation then sends it to \((15+t)-8=7+t\).  The whiskering is
	enabled by a distinguished ridge satisfying the eigenvector condition
	proved below.  Thus the two seeds cover every \(r\ge5\).
	Both lie on \(\{0,1,\ldots,6\}\).
Write \(ijk\) for the set \(\{i,j,k\}\).  On the ground set
	\(\{0,1,\ldots,6\}\), define
	\begin{align*}
		\calF_{16}=\{&012,013,015,026,123,124,125,126,\\
		&134,135,136,145,146,156,236,256\},\\
		\calF_{15}=\{&023,025,026,035,036,123,125,126,135,136,\\
		&234,235,236,256,345\}.
	\end{align*}
	Their degree sequences, in vertex order, are
	\[
	(4,13,8,6,4,6,7),
	\qquad
	(5,5,10,10,2,7,6).
	\]
To explicitly determine the spectral defects of the two seeds, we first record their characteristic polynomials.
\begin{lemma}\label{lem:seed-polynomials}
		Put
		\begin{align*}
			p_{16}(x)&=x^5-19x^4+133x^3-413x^2+527x-175,\\
			p_{15}(x)&=x^5-16x^4+91x^3-224x^2+233x-84.
		\end{align*}
		Then
		\begin{align*}
			\det(xI-M_{\calF_{16}})
			&=x^2(x-1)^5(x-3)(x-7)^3p_{16}(x),\\
			\det(xI-M_{\calF_{15}})
			&=x^4(x-2)(x-4)^2(x-6)^2(x-7)p_{15}(x)
		\end{align*}
		and \(p_{15}(x)=-p_{16}(7-x)\).
	\end{lemma}
	
	\begin{proof}
		Order the facets as displayed above and retain the vertex order
		\(0<1<\cdots<6\).  For either
		\(\calF\in\{\calF_{15},\calF_{16}\}\) and
		\(\sigma,\tau\in\calF\), the entry
		\((M_{\calF})_{\sigma,\tau}=\ip{\bd e_\sigma}{\bd e_\tau}\) is \(3\) on the
		diagonal and is \(0\) or \(\pm1\) off the diagonal, with a nonzero entry
		exactly when \(\sigma\) and \(\tau\) share an edge.  This rule determines
		the two integer matrices uniquely; they are recorded in
		Appendix~\ref{app:seed-verification} for direct verification.
		Bareiss fraction-free elimination of \(xI-M_{\calF}\), entirely in
		\(\mathbb Z[x]\), gives
		the desired result.  Substitution gives
		\(p_{15}(x)=-p_{16}(7-x)\).
	\end{proof}
	The following root isolation establishes the necessary bounds for the eigenvalues of the sixteen-facet seed.
	\begin{lemma}\label{lem:seed-roots}
		The roots of \(p_{16}\), in increasing order, satisfy
		\[
		\theta_1\in(0,1),\quad
		\theta_2\in(2,3),\quad
		\theta_3\in(4,5),\quad
		\theta_4\in(5,6),\quad
		\theta_5\in(6,7),
		\]
		and \(\theta_4+\theta_5>12\).
	\end{lemma}
	
	\begin{proof}
		Direct substitution gives
		\[
		\begin{array}{c|rrrrrrrr}
			x&0&1&2&3&4&5&6&7\\ \hline
			p_{16}(x)&-175&54&19&-16&-3&10&-1&84.
		\end{array}
		\]
		The intermediate value theorem produces a root in each of
		\((0,1)\), \((2,3)\), \((4,5)\), \((5,6)\), and \((6,7)\).
		These five intervals are disjoint, and \(p_{16}\) has degree five, so
		each interval contains exactly one root and there are no others.
		For the final assertion, set
		\[
		a=\theta_5-6\in(0,1),
		\qquad
		g(t)=p_{16}(6+t),
		\qquad
		h(t)=t^4+37t^2-1.
		\]
		The explicit polynomial $g(t)=t^5+11t^4+37t^3+37t^2-t-1$
		satisfies
		\[
		g(t)=(1+t)h(t)+10t^4,
		\qquad
		g(t)-g(-t)=2t h(t).
		\]
		
		Since \(g(a)=0\), the first identity gives
		$h(a)=-\frac{10a^4}{1+a}<0.$
		The second identity therefore yields
		$g(-a)=-2a h(a)=\frac{20a^5}{1+a}>0.$
		Because \(g(0)=-1<0\), the intermediate value theorem gives a root of
		\(g\) in \((-a,0)\).  The first part shows that the unique root of
		\(g\) in \((-1,0)\) is \(\theta_4-6\).  Hence
		$\theta_4-6>-a=6-\theta_5,$
		which is equivalent to \(\theta_4+\theta_5>12\).
	\end{proof}
	Set
$
d_*=\theta_4+\theta_5-12>0.
$
	Numerically, \(d_*\approx1.9634\times10^{-4}\), which indicates the
	small scale of the violation; its positivity follows exactly from
	Lemma~\ref{lem:seed-roots}, not from floating-point computation.
Combining the characteristic polynomials with the root bounds,
we explicitly determine the partial sums at the critical indices.
	
	\begin{proposition}\label{prop:seed-defects}
		The following identities hold:
		\begin{align*}
			s_5(\calF_{16})&=D_5(\calF_{16})+d_*=33+d_*,\\
			s_8(\calF_{15})&=D_8(\calF_{15})+d_*=41+d_*,\\
			s_9(\calF_{15})&=D_9(\calF_{15})+d_*=43+d_*.
		\end{align*}
	\end{proposition}
	
	\begin{proof}
		By Lemmas~\ref{lem:seed-polynomials} and \ref{lem:seed-roots}, the five
		largest eigenvalues of \(M_{\calF_{16}}\) are
		\(7,7,7,\theta_5,\theta_4\).  Hence
		\[
		s_5(\calF_{16})=21+\theta_4+\theta_5=33+d_*.
		\]
		The displayed degree sequence gives \(D_5(\calF_{16})=33\).
		The roots of \(p_{15}\) are \(7-\theta_i\).  The eight largest
		eigenvalues of \(M_{\calF_{15}}\) are
		\[
		7,\ 7-\theta_1,\ 6,\ 6,\ 7-\theta_2,\ 4,\ 4,\ 7-\theta_3.
		\]
		Using \(\sum_i\theta_i=19\), their sum is \(41+d_*\).  The next
		eigenvalue is \(2\), so \(s_9(\calF_{15})=43+d_*\).  The degree
		sequence gives \(D_8(\calF_{15})=41\) and \(D_9(\calF_{15})=43\).
	\end{proof}
	

	The two seed gaps will be propagated by ridge-whiskering and set-complement
	duality.  We develop both operations in the generality needed below.
For every face \(\rho\subseteq V\), define its \(F\)-degree by
$d_F(\rho)=|\{\sigma\in F:\rho\subseteq\sigma\}|.$
For a vertex \(\rho=\{v\}\), this agrees with the previously defined
vertex degree \(d_F(v)\).
Let \(F\subseteq\binom Vq\) and \(\eta\in\binom V{q-1}\).  Using the
	oriented incidence coefficient defined in Section~\ref{sec:notation}, put
	\[
	\mathbf y_\eta
	=\bigl(\varepsilon(\eta,\sigma)\bigr)_{\sigma\in F}
	\in\R^F.
	\]
	Thus \(\norm{\mathbf y_\eta}^2=d_F(\eta)\).  For \(t\ge0\), choose pairwise
	distinct fresh vertices \(w_1,\ldots,w_t\notin V\) and set
	\[
	F(t)=F\cup\{\eta\cup\{w_i\}:1\le i\le t\}.
	\]
	
	\begin{lemma}\label{lem:ridge-transfer}
		Assume \(d_F(\eta)>0\) and $
		M_F\mathbf y_\eta=(d_F(\eta)+q-1)\mathbf y_\eta.$
		Then
		\[
		\det(xI-M_{F(t)})
		=\det(xI-M_F)
		\frac{(x-(d_F(\eta)+q-1+t))(x-(q-1))^t}
		{x-(d_F(\eta)+q-1)}.
		\]
	\end{lemma}
\begin{proof}
Put \(d=d_F(\eta)\). The case \(t=0\) is immediate, so assume
\(t\geq1\). Orient each new facet \(\eta\cup\{w_i\}\) so that its
incidence coefficient at \(\eta\) is \(+1\); changing facet orientations
does not change the characteristic polynomial. Since the \(w_i\) are
fresh, an old facet and a new facet can share only \(\eta\), while two
distinct new facets share exactly \(\eta\). Thus
\[
M_{F(t)}
=
\begin{pmatrix}
M_F&\mathbf y_\eta\mathbf 1_t^{\mathsf T}\\
\mathbf 1_t\mathbf y_\eta^{\mathsf T}
&(q-1)I_t+\mathbf 1_t\mathbf 1_t^{\mathsf T}
\end{pmatrix}.
\]
If \(\mathbf z\perp\mathbf y_\eta\), then symmetry of \(M_F\) and the
hypothesis \(M_F\mathbf y_\eta=(d+q-1)\mathbf y_\eta\) give
\(\langle M_F\mathbf z,\mathbf y_\eta\rangle=0\). Hence the displayed
matrix sends \((\mathbf z,0)\) to \((M_F\mathbf z,0)\). Similarly, if
\(\mathbf u\perp\mathbf 1_t\), it sends \((0,\mathbf u)\) to
\((0,(q-1)\mathbf u)\). These two subspaces are therefore invariant.
Their orthogonal complement is two-dimensional, with orthonormal basis
$\left(\frac{\mathbf y_\eta}{\sqrt d},0\right)$ and $
\left(0,\frac{\mathbf 1_t}{\sqrt t}\right).$
In this basis, the diagonal entries are \(d+q-1\) and \(q-1+t\), and
the off-diagonal entry is \(\sqrt{td}\). Consequently,
\[
M_{F(t)}
\sim
\left.M_F\right|_{\mathbf y_\eta^\perp}
\oplus(q-1)I_{t-1}
\oplus
\begin{pmatrix}
d+q-1&\sqrt{td}\\
\sqrt{td}&q-1+t
\end{pmatrix},
\qquad
M_F\sim
\left.M_F\right|_{\mathbf y_\eta^\perp}\oplus(d+q-1),
\]
where \(\sim\) denotes orthogonal similarity.
The characteristic polynomial of the two-dimensional block is $
(x-(q-1))(x-(d+q-1+t)).$
Taking characteristic polynomials in the two orthogonal decompositions
therefore gives
\[
\det(xI-M_{F(t)})
=
\det(xI-M_F)
\frac{(x-(d+q-1+t))(x-(q-1))^t}
{x-(d+q-1)},
\]
as required.
\end{proof}

	

	
	Fix the ambient vertex set \(V\).  Let \(|V|=n\), let \(1\le q<n\), let
	\(F\subseteq\binom Vq\), and define its set-complement family by
	\[
	F^\star=\{V\setminus\sigma:\sigma\in F\}
	\subseteq\binom V{n-q}.
	\]
	This operation depends on the specified ambient set \(V\).  In particular,
	isolated vertices may be deleted when studying \(s_r(F)\) and \(D_r(F)\),
	but they must not be deleted silently before forming \(F^\star\).
	The following is Proposition~4.2 of Duval and Reiner
	\cite{DuvalReiner2002}, restated in our notation.
	
	\begin{lemma}\label{lem:hodge-star}
		Write \(V=\{v_1<\cdots<v_n\}\), and index \(M_{F^\star}\) by \(F\)
		via \(\sigma\leftrightarrow V\setminus\sigma\).  Then
		$M_{F^\star}=S(nI-M_F)S,$ where $
		S=\diag\left((-1)^{\sum_{v_i\in\sigma}i}:\sigma\in F\right).$
	\end{lemma}
	
By exploiting the complementarity of the spectra,
we obtain the following transfer identity for the spectral defect.
\begin{lemma}\label{lem:defect-transfer}
Let \(f=|F|\). If \(1\le r<f\), then
$
\Delta_{f-r}(F^\star)=\Delta_r(F).
$
\end{lemma}

\begin{proof}
Let \(f=|F|\) and \(\lambda_1\geq\cdots\geq\lambda_f\) be the eigenvalues of \(M_F\),
including zeros. By Lemma~\ref{lem:hodge-star}, the eigenvalues of
\(M_{F^\star}\), in decreasing order, are
\(n-\lambda_f,\ldots,n-\lambda_1\). Since
\(\tr(M_F)=qf\), it follows that
\begin{equation}\label{a}
s_{f-r}(F^\star)
=n(f-r)-\sum_{i=r+1}^f\lambda_i
=(n-q)f-nr+s_r(F).
\end{equation}
For every \(v\in V\), we have
\(d_{F^\star}(v)=f-d_F(v)\). Moreover, for \(0\leq d\leq f\),
\[
\min\{f-d,f-r\}=f-r-(d-r)_+,
\qquad
\min\{d,r\}=d-(d-r)_+.
\]
The second identity and \(\sum_{v\in V}d_F(v)=qf\) give
\[
\sum_{v\in V}(d_F(v)-r)_+=qf-D_r(F).
\]
Applying the first identity with \(d=d_F(v)\), we obtain
\begin{equation}\label{b}
D_{f-r}(F^\star)
=n(f-r)-\sum_{v\in V}(d_F(v)-r)_+
=(n-q)f-nr+D_r(F).
\end{equation}
Equalities \eqref{a}-\eqref{b} give
\(\Delta_{f-r}(F^\star)=\Delta_r(F)\), as required.
\end{proof}
	For \(\calF_{15}\), take the ridge \(\eta=\{2,3\}\).  In the displayed
	facet order, the signed ridge-incidence vector defined above is
	\[
	\mathbf y_\eta=(1,0,0,0,0,1,0,0,0,0,1,1,1,0,0)^{\mathsf T}.
	\]
	Exact multiplication gives
	\begin{equation}\label{eq:F15-ridge-vector}
		\norm{\mathbf y_\eta}^2=5,
		\qquad
		M_{\calF_{15}}\mathbf y_\eta=7\mathbf y_\eta.
	\end{equation}
	For \(t\geq0\), let \(w_1,\ldots,w_t\) be pairwise distinct vertices
outside \(\{0,1,\ldots,6\}\), and define
\[
\calF_{15}(t)
=
\calF_{15}\cup\{\{2,3,w_i\}:1\leq i\leq t\}.
\]
	
	\begin{proposition}\label{prop:whiskered-seed}
		For every \(t\ge0\),
		\begin{align}
			\det(xI-M_{\calF_{15}(t)})
			&=x^4(x-2)^{t+1}(x-4)^2(x-6)^2(x-t-7)p_{15}(x),
			\label{eq:whiskered-charpoly}\\
			\Delta_8(\calF_{15}(t))&=d_*.\label{eq:whiskered-defect}
		\end{align}
	\end{proposition}
	
	\begin{proof}
		Equation \eqref{eq:F15-ridge-vector} satisfies
		Lemma~\ref{lem:ridge-transfer} with \(q=3\) and \(d_{\calF_{15}}(\eta)=5\), proving
		\eqref{eq:whiskered-charpoly}.  The degrees are
		\[
		(5,5,10+t,10+t,2,7,6,
		\underbrace{1,\ldots,1}_{t\text{ times}}),
		\]
		so \(D_8(\calF_{15}(t))=41+t\).  Since \(2<7-\theta_3<3\), the eight
		largest eigenvalues obtained from \eqref{eq:whiskered-charpoly} are
		\[
		t+7,\ 7-\theta_1,\ 6,\ 6,\ 7-\theta_2,\ 4,\ 4,\ 7-\theta_3.
		\]
		They sum to \(41+t+d_*\), proving
		\eqref{eq:whiskered-defect}.
	\end{proof}
	
These properties of the seed families provide all the necessary ingredients to establish Theorem~\ref{thm:counterexamples}.
	\begin{proof}[\textup{\textbf{Proof of Theorem~\ref{thm:counterexamples}}}]
		At \(r=5\), take \(\calF_{16}\).  By
		Proposition~\ref{prop:seed-defects},
		$\Delta_5(\calF_{16})=d_*>0.$
		Every entry of its displayed degree sequence is strictly less than
		\(|\calF_{16}|=16\), so no vertex belongs to every facet.  Hence
		\(\calF_{16}\) has empty total facet intersection.
		For \(t\ge0\), put
		\[
			V_t=\{0,1,\ldots,6,w_1,\ldots,w_t\},
			\qquad
			\mathcal G_t
			=\{V_t\setminus\sigma:\sigma\in\calF_{15}(t)\},
		\]
		where complementation is taken relative to \(V_t\). Thus
		\(\mathcal G_t\) is \((t+4)\)-uniform and has \(15+t\) facets.
		Every vertex of \(V_t\) occurs in some member of
		\(\calF_{15}(t)\), and hence is omitted from some member of
		\(\mathcal G_t\). Therefore \(\mathcal G_t\) has empty total facet
		intersection.
		At \(r=6\), take \(\mathcal G_0\). Lemma~\ref{lem:defect-transfer}
		and Proposition~\ref{prop:seed-defects} give
		\[
			\Delta_6(\mathcal G_0)
			=\Delta_9(\calF_{15})=d_*>0.
		\]
		For \(r\ge7\), put \(t=r-7\). Since
		\(|\calF_{15}(t)|=15+t\), Lemma~\ref{lem:defect-transfer} and
		\eqref{eq:whiskered-defect} yield
		\[
			\Delta_r(\mathcal G_t)
			=\Delta_{7+t}(\mathcal G_t)
			=\Delta_8(\calF_{15}(t))
			=d_*>0.
		\]
		This proves the assertion for every \(r\ge5\).
		It remains to consider the uniformity.  For \(q=3\), use
		\(\calF_{16}\).  For \(q\ge4\), put \(t=q-4\).  Then
		\(\mathcal G_t\) is \(q\)-uniform, and the preceding display gives
		a strict counterexample at index \(7+t=q+3\), again with empty total
		facet intersection.
	\end{proof}
	
	\section{The Second Partial-Sum Inequality and Its Equality Cases}
	\label{sec:core-inequality}
	We now shift our focus from the counterexamples to the universal majorization bound at the second index.
	For the index \(r=2\), recall the pendant and core sets \(P(F)\) and
	\(Q(F)\) from the Introduction.
	Then  \(V=P(F)\cup Q(F)\). Since \(\min\{d_F(v),2\}\) is \(1\) on \(P(F)\) and \(2\) on \(Q(F)\), we have
	\begin{equation}\label{eq:R2-pq}
		D_2(F)
		=
		\sum_{v\in V}\min\{d_F(v),2\}
		=
		|P(F)|+2|Q(F)|.
	\end{equation}
Let
	\[
		\widehat F=F\cup\binom{Q(F)}q.
	\]
	When \(|Q(F)|<q\), the completion adds no facets.
	
	\begin{lemma}\label{lem:completion}
		With the notation above,
		\[
		P(\widehat F)=P(F),
		\qquad
		Q(\widehat F)=Q(F),
		\qquad
		D_2(\widehat F)=D_2(F).
		\]
		Moreover, for every \(r\ge1\),
		\(s_r(F)\le s_r(\widehat F)\).
	\end{lemma}
	
	\begin{proof}
		Since \(\widehat F\setminus F\subseteq\binom{Q(F)}q\), no added facet contains a
		vertex outside \(Q(F)\). Hence
		\(d_{\widehat F}(v)=d_F(v)\) if \(v\notin Q(F)\), and
		\(d_{\widehat F}(v)\ge d_F(v)\ge2\) if \(v\in Q(F)\).
		It follows that \(P(\widehat F)=P(F)\) and
		\(Q(\widehat F)=Q(F)\).
		Hence, by \eqref{eq:R2-pq},
		\(D_2(\widehat F)=D_2(F)\).
		Let \(\widetilde B_F\) and \(\widetilde B_{\widehat F}\) be the boundary
		matrices with rows indexed by \(\binom V{q-1}\) and columns indexed by
		\(F\) and \(\widehat F\), respectively.  Thus \(\widetilde B_F\)
		is obtained from \(B_F\) by adjoining zero rows, so
		\(\widetilde B_F\widetilde B_F^{\mathsf T}\) has the same nonzero
		eigenvalues as \(L_F^+\). Let \(C\) be the
		matrix whose columns are the coordinate vectors of \(\bd e_{\sigma}\)
		for \(\sigma\in\widehat F\setminus F\). Then
		\(\widetilde B_{\widehat F}
		=\begin{pmatrix}\widetilde B_F&C\end{pmatrix}\). Hence
		\[
		\widetilde B_{\widehat F}\widetilde B_{\widehat F}^{\mathsf T}
		=
		\widetilde B_F\widetilde B_F^{\mathsf T}
		+
		CC^{\mathsf T}.
		\]
		Since \(CC^{\mathsf T}\) is positive
		semidefinite, fix \(r\ge1\)
		and, if necessary, adjoin the same number of zero rows and columns to
		both matrices so that their common size is at least \(r\).  This does
		not change either Ky Fan sum under the appended-zero convention, and
		the difference remains positive semidefinite.  Lemma~\ref{lem:kyfan-monotone}
		therefore gives \(s_r(F)\le s_r(\widehat F)\).  Since \(r\) was arbitrary,
		the inequality holds for every \(r\ge1\).
	\end{proof}
	
	We call the vertices in \(P(F)\) \emph{pendant vertices} and the vertices in
	\(Q(F)\) \emph{core vertices}. A facet containing exactly one pendant
	vertex is \emph{one-pendant}, and a facet containing at least two pendant
	vertices is \emph{heavy}.
	In any \(q\)-family, every facet is naturally either a core facet (contained entirely in \(Q(F)\)), a one-pendant facet \(\rho\cup\{z\}\) with \(\rho\in\binom{Q(F)}{q-1}\) and \(z\in P(F)\), or a heavy facet.
	
	\begin{lemma}\label{lem:star}
		Let \(\rho\) be a \((q-1)\)-set, let \(Z\) be a finite set disjoint
		from \(\rho\), and put
		\(\mathcal S(\rho,Z)=\{\rho\cup\{z\}:z\in Z\}\).
		If \(|Z|\ge1\), then the eigenvalues of
		\(M_{\mathcal S(\rho,Z)}\) are \(|Z|+q-1\), and \(q-1\) repeated
		\(|Z|-1\) times.
	\end{lemma}
	
	\begin{proof}
		Write \(b=|Z|\) and \(\sigma_z=\rho\cup\{z\}\) for \(z\in Z\). The
		diagonal entries of \(M_{\mathcal S(\rho,Z)}\) are \(q\). For
		\(z\ne w\), the facets \(\sigma_z,\sigma_w\) have the unique common
		\((q-1)\)-face \(\rho\).
		Thus
		\((M_{\mathcal S(\rho,Z)})_{zw}
		=\varepsilon(\rho,\sigma_z)\varepsilon(\rho,\sigma_w)\).
		Rescaling each basis vector \(e_{\sigma_z}\) by the sign
		\(\varepsilon(\rho,\sigma_z)\) yields an orthogonally similar matrix
		\[
		(q-1)I_b+\one_b\one_b^{\mathsf T}.
		\]
		The all-ones vector \(\one_b\) is an eigenvector of this matrix with eigenvalue
		\(b+q-1\), and its orthogonal complement \(\one_b^\perp\) is the eigenspace
		for the eigenvalue \(q-1\). Hence the eigenvalues of
		\(M_{\mathcal S(\rho,Z)}\) are \(b+q-1\) with multiplicity one, and \(q-1\) with
		multiplicity \(b-1\).
	\end{proof}
To bound the eigenvalues of $M_F$, we first identify the explicit eigenvectors generated by edges with at least two pendants.
\begin{lemma}\label{lem:heavy-isolated}
	Let \(F\subseteq\binom Vq\), and let \(\sigma\in F\) contain at least two pendant vertices. Then \(e_\sigma\) is an eigenvector of \(M_F\) with eigenvalue \(q\).
\end{lemma}

\begin{proof}
	Let \(\tau\in F\setminus\{\sigma\}\). By definition, the pendant vertices in \(\sigma\) belong to no other facet in \(F\). Since \(\sigma\) contains at least two such vertices, it follows that \(|\sigma\cap\tau|\le q-2\). Thus \(\sigma\) and \(\tau\) share no \((q-1)\)-face, and hence \((M_F)_{\sigma,\tau}=\ip{\bd e_\sigma}{\bd e_\tau}=0\). Furthermore, \((M_F)_{\sigma,\sigma}=\norm{\bd e_\sigma}^2=q\). Therefore, \(M_Fe_\sigma=qe_\sigma\).
\end{proof}
	
\begin{lemma}\label{lem:leaf-differences}
Fix \(\rho\in\binom{Q(F)}{q-1}\), and let
\(\sigma_i=\rho\cup\{z_i\}\), \(1\leq i\leq b\), be all one-pendant
facets with core \(\rho\). Then
\[
\left\{
\sum_{i=1}^b
a_i\varepsilon(\rho,\sigma_i)e_{\sigma_i}:
\sum_{i=1}^b a_i=0
\right\}
\]
is contained in the \((q-1)\)-eigenspace of \(M_F\).
\end{lemma}

\begin{proof}
The assertion is trivial if \(b=0\), so assume \(b\geq1\). Write
\(\varepsilon_i=\varepsilon(\rho,\sigma_i)\), and let
$\mathbf u=\sum_{i=1}^b a_i\varepsilon_i e_{\sigma_i},$
where $\sum_{i=1}^b a_i=0.$
In the rescaled basis \(\varepsilon_i e_{\sigma_i}\), the principal
submatrix indexed by \(\sigma_1,\ldots,\sigma_b\) is
\((q-1)I_b+\mathbf 1_b\mathbf 1_b^{\mathsf T}\): its diagonal entries
are \(q\), and its off-diagonal entries are \(1\), since two distinct
\(\sigma_i\) share exactly the ridge \(\rho\). It therefore sends the
coefficient vector \((a_1,\ldots,a_b)^{\mathsf T}\) to
\((q-1)(a_1,\ldots,a_b)^{\mathsf T}\).

It remains to check the coordinates outside
\(\{\sigma_1,\ldots,\sigma_b\}\). Let
\(\tau\in F\setminus\{\sigma_1,\ldots,\sigma_b\}\). Any ridge shared by
\(\tau\) and \(\sigma_i\) cannot contain \(z_i\), because \(z_i\) occurs
only in \(\sigma_i\); hence the shared ridge, if one exists, must be
\(\rho\). Thus either all relevant Gram entries vanish, or
\[
(M_F\mathbf u)_\tau
=
\sum_{i=1}^b
a_i\varepsilon_i
\varepsilon(\rho,\tau)\varepsilon_i
=
\varepsilon(\rho,\tau)\sum_{i=1}^b a_i
=0.
\]
Hence \(M_F\mathbf u\) has no coordinates outside
\(\{\sigma_1,\ldots,\sigma_b\}\), while its coordinates inside this set
are those of \((q-1)\mathbf u\). Therefore
\(M_F\mathbf u=(q-1)\mathbf u\), as required.
\end{proof}
	
Let
\[
Q=Q(F),\qquad P=P(F),\qquad m=|Q|,\qquad p=|P|.
\]
Define
\[
\calF_Q=\binom Qq, \qquad  \calR_Q=\binom Q{q-1}.
\]
Recall the
signed incidence map \(\bd_Q:\R^{\calF_Q}\to\R^{\calR_Q}\) from the
Introduction, given explicitly by
\[
\bd_Q e_\sigma
=
\sum_{j=1}^q(-1)^{j-1}e_{\sigma\setminus\{v_j\}},
\qquad
\sigma=\{v_1<\cdots<v_q\}.
\]
Furthermore, put
\[
\scrH_Q=\im(\bd_Q^{\mathsf T}),
\qquad
\mathbf y_\rho=\bd_Q^{\mathsf T}e_\rho
\quad(\rho\in\calR_Q).
\]
The complete-core matrix is then
\(M_{\calF_Q}=\bd_Q^{\mathsf T}\bd_Q\).	
Next, we evaluate the boundary operator and the specific support structures within a complete simplex.
\begin{lemma}\label{lem:complete-simplex}
	Let \(m=|Q|\ge q\). Then
	\[ \dim(\im\bd_Q^{\mathsf T})=\binom{m-1}{q-1}, \qquad
	\bd_Q^{\mathsf T}\bd_Q=
	\begin{cases}
		mI,&\text{on }\im\bd_Q^{\mathsf T},\\
		0,&\text{on }\Ker\bd_Q.
	\end{cases}
	\]
	For \(\rho\in\calR_Q\), the vector \(\mathbf y_\rho\) satisfies
	\begin{align*}
		\supp(\mathbf y_\rho)&=\{\rho\cup\{w\}:w\in Q\setminus\rho\},\\
		\norm{\mathbf y_\rho}^2&=m-q+1
	\end{align*}
	and, for distinct \(\rho,\gamma\in\calR_Q\),
	\(\ip{\mathbf y_\rho}{\mathbf y_\gamma}=0\) if and only if
	\(|\rho\cap\gamma|\le q-3\).
\end{lemma}

\begin{proof}
	We first establish the rank. Put \(h=\binom{m-1}{q-1}\), and fix \(v\in Q\). Among the columns of
	\(\bd_Q\) indexed by \(\sigma=\rho\cup\{v\}\),
	\(\rho\in\binom{Q\setminus\{v\}}{q-1}\), and the rows indexed by those
	same \(\rho\)'s, the resulting submatrix is a signed identity matrix.
	Hence
	\[
	\rank\bd_Q\ge\binom{m-1}{q-1}=h.
	\]
	Let \(\bd_Q^-:\R^{\calR_Q}\to\R^{\binom Q{q-2}}\) be the lower
	boundary. The analogous submatrix of \(\bd_Q^-\), using columns
	\(\eta\cup\{v\}\) and rows
	\(\eta\in\binom{Q\setminus\{v\}}{q-2}\), is also a signed identity.
	Thus \(\rank\bd_Q^-\ge\binom{m-1}{q-2}\). Since
	\(\bd_Q^-\bd_Q=0\),
	\[
	\rank\bd_Q
	\le\dim\Ker\bd_Q^-
	\le\binom m{q-1}-\binom{m-1}{q-2}
	=\binom{m-1}{q-1}=h.
	\]
	Consequently, \(\dim(\im\bd_Q^{\mathsf T}) = \rank\bd_Q = h\).
	
	We next prove the complete-simplex identity with signs. If
	\(\rho,\gamma\in\binom Q{q-1}\) are distinct, both off-diagonal entries
	\((\bd_Q\bd_Q^{\mathsf T})_{\rho,\gamma}\) and
	\(((\bd_Q^-)^{\mathsf T}\bd_Q^-)_{\rho,\gamma}\) vanish unless
	\(|\rho\cap\gamma|=q-2\). When \(|\rho\cap\gamma|=q-2\), let
	\(\eta=\rho\cap\gamma\) and \(\sigma=\rho\cup\gamma\). The coefficient of
	\(e_\eta\) in \(\bd^2e_\sigma=0\) is
	\[
	\varepsilon(\rho,\sigma)\varepsilon(\eta,\rho)
	+\varepsilon(\gamma,\sigma)\varepsilon(\eta,\gamma)=0.
	\]
	Multiplying this identity by the sign
	\(\varepsilon(\gamma,\sigma)\varepsilon(\eta,\rho)\) and using the fact that
	every incidence coefficient squares to one yields
	\[
	\varepsilon(\rho,\sigma)\varepsilon(\gamma,\sigma)
	+\varepsilon(\eta,\rho)\varepsilon(\eta,\gamma)=0.
	\]
	The two terms in the last display are precisely the off-diagonal entries
	of \(\bd_Q\bd_Q^{\mathsf T}\) and
	\((\bd_Q^-)^{\mathsf T}\bd_Q^-\), so those entries cancel.
	On the diagonal indexed by \(\rho\), the two contributions are
	\(m-q+1\), the number of \(q\)-sets containing \(\rho\), and \(q-1\),
	the number of codimension-one faces of \(\rho\). Therefore
	\[
	\bd_Q\bd_Q^{\mathsf T}+(\bd_Q^-)^{\mathsf T}\bd_Q^-=mI.
	\]
	The orthogonal decomposition
	$\R^{\calF_Q}=\im\bd_Q^{\mathsf T}\oplus\Ker\bd_Q$
	follows from \((\im\bd_Q^{\mathsf T})^\perp=\Ker\bd_Q\). The
	operator \(\bd_Q^{\mathsf T}\bd_Q\) is zero on \(\Ker\bd_Q\).
	If \(\mathbf x=\bd_Q^{\mathsf T}\mathbf z\), then the identity above and
	\(\bd_Q^-\bd_Q=0\) give
	\begin{align*}
		\bd_Q^{\mathsf T}\bd_Q\mathbf x=\bd_Q^{\mathsf T}\bd_Q\bd_Q^{\mathsf T}\mathbf z=\bd_Q^{\mathsf T}
		\bigl(mI-(\bd_Q^-)^{\mathsf T}\bd_Q^-\bigr)\mathbf z
		=m\mathbf x.
	\end{align*}
	This proves the asserted action.
	Finally, \(\mathbf y_\rho=\bd_Q^{\mathsf T}e_\rho\) has a coordinate of
	\(\pm1\) for every \(q\)-set containing \(\rho\), and is zero otherwise. Hence
	\[
	\supp(\mathbf y_\rho)=\{\rho\cup\{w\}:w\in Q\setminus\rho\},
	\qquad \norm{\mathbf y_\rho}^2=m-q+1.
	\]
	For distinct \(\rho,\gamma\), their inner product has a non-zero contribution precisely
	when a \(q\)-set contains both, which occurs precisely when
	\(|\rho\cap\gamma|=q-2\); the unique contribution is then \(\pm1\). Thus
	\(\ip{\mathbf y_\rho}{\mathbf y_\gamma}=0\) exactly when
	\(|\rho\cap\gamma|\le q-3\).
\end{proof}
	
Let \(\calC\) be the set of ridges that are contained in at least one
one-pendant facet:
\[
\calC=\{\rho\in\calR_Q:\rho\cup\{z\}\in F\text{ for some }z\in P\}.
\]
For each \(\rho\in\calC\), let \(b_\rho=|\{z\in P:\rho\cup\{z\}\in F\}|\) denote the
number of pendant vertices associated with \(\rho\), and put
\(p_1=\sum_{\rho\in\calC}b_\rho\). Since every pendant vertex in \(P\) is
contained in a unique facet, it follows that \(p_1\le p\). Define the linear
operator
$$
Y:\R^{\calC}\to\scrH_Q, \qquad Ye_\rho=\sqrt{b_\rho}\mathbf y_\rho\  (\rho\in\calC).
$$
We define the matrix \(T_F\) as the block operator acting on \(\scrH_Q\oplus\R^{\calC}\) given by
\begin{equation*}
	T_F=
	\begin{pmatrix}
		0 & Y \\
		Y^{\mathsf T} & \diag(b_\rho:\rho\in\calC)-(m-q+1)I
	\end{pmatrix}.
\end{equation*}

\begin{lemma}\label{lem:core-excess} Assume $\binom{Q}q\subseteq F$ and $m\ge q$. Every eigenvalue of $M_F$ strictly exceeding $m$ is $m+\mu$ for a positive eigenvalue $\mu$ of $T_F$, with the same multiplicity. Hence $$ \sum_{\lambda\in\Spec(M_F)}(\lambda-m)_+ =\tr((T_F)_+)\le p_1\le p, $$ and, for every $r\ge1$, $$ s_r(F) \le rm+\tr((T_F)_+) \le rm+p_1 \le rm+p. $$ \end{lemma}

\begin{proof}
For each \(\rho\in\calC\), define
\[
\boldsymbol{\ell}_\rho
=
b_\rho^{-1/2}
\sum_{\substack{z\in P\\ \rho\cup\{z\}\in F}}
\varepsilon(\rho,\rho\cup\{z\})e_{\rho\cup\{z\}}.
\]
Each \(\boldsymbol{\ell}_\rho\) has norm one. Since every one-pendant
facet has a unique core, these vectors have disjoint supports for
distinct \(\rho\), and hence are orthonormal.
Every facet is either a core facet, a one-pendant facet, or a heavy
facet. The assumption \(\binom Qq\subseteq F\) implies that the core
facets are exactly \(\calF_Q=\binom Qq\). Let \(g\) be the number of
heavy facets. Lemma~\ref{lem:heavy-isolated} gives the reducing block
\(qI_g\). For each \(\rho\in\calC\),
Lemma~\ref{lem:leaf-differences} gives a
\((b_\rho-1)\)-dimensional eigenspace with eigenvalue \(q-1\).
Since \(M_F\) is symmetric, these eigenspaces are reducing and together
give the block \((q-1)I_{p_1-|\calC|}\). Their orthogonal complement,
together with that of the heavy-facet space, is
\[
\mathbb R^{\calF_Q}
\oplus
\Span\{\boldsymbol{\ell}_\rho:\rho\in\calC\}.
\]

If \(\rho\neq\gamma\), a one-pendant facet with core \(\rho\) and one
with core \(\gamma\) intersect in \(\rho\cap\gamma\), which has size at
most \(q-2\). They therefore share no ridge, so
\(\langle M_F\boldsymbol{\ell}_\rho,
\boldsymbol{\ell}_\gamma\rangle=0\).
Lemma~\ref{lem:star} also gives
\(\langle M_F\boldsymbol{\ell}_\rho,
\boldsymbol{\ell}_\rho\rangle=b_\rho+q-1\).
For \(\sigma\in\calF_Q\), direct calculation gives
\[
\langle M_Fe_\sigma,\boldsymbol{\ell}_\rho\rangle
=
b_\rho^{-1/2}
\sum_{\substack{z\in P\\ \rho\cup\{z\}\in F}}
\varepsilon(\rho,\rho\cup\{z\})
\langle\partial e_\sigma,\partial e_{\rho\cup\{z\}}\rangle
=
\sqrt{b_\rho}\,\varepsilon(\rho,\sigma).
\]
Thus the coupling between the core space and the leaf-average space is
\(Y\), and the latter has diagonal block
\(\diag(b_\rho+q-1:\rho\in\calC)\).
By Lemma~\ref{lem:complete-simplex},
\(\mathbb R^{\calF_Q}=\scrH_Q\oplus\Ker\partial_Q\),
with \(M_{\calF_Q}=mI\) on \(\scrH_Q\) and \(M_{\calF_Q}=0\) on
\(\Ker\partial_Q\). If \(\mathbf x\in\Ker\partial_Q\), then
\(\mathbf x\perp\scrH_Q\); since \(\im Y\subseteq\scrH_Q\), this gives
\(Y^{\mathsf T}\mathbf x=0\). Hence, using the definition of \(T_F\),
we obtain the orthogonal decomposition
\[
M_F\simeq
qI_g
\oplus(q-1)I_{p_1-|\calC|}
\oplus0_{\binom{m-1}{q}}
\oplus
\left(
mI_{\binom{m-1}{q-1}+|\calC|}+T_F
\right).
\]
Because \(m\geq q\), every eigenvalue of the first three blocks is at
most \(m\). The eigenvalues of the last block are \(m+\mu\), where
\(\mu\) runs through the eigenvalues of \(T_F\), with the same
multiplicities. Therefore
\[
\sum_{\lambda\in\Spec(M_F)}(\lambda-m)_+
=
\tr((T_F)_+).
\]
If \(\calC=\varnothing\), then \(T_F=0\). Otherwise,
Lemma \ref{lem:complete-simplex} permits us to apply
Lemma~\ref{lem:positive-excess-rigidity} with
\(a=m-q+1\), \(\mathbf z_\rho=\mathbf y_\rho\),
\(b_i=b_\rho\), and \(X=Y\). In either case,
\[
\tr((T_F)_+)
\leq\sum_{\rho\in\calC}b_\rho
=p_1
\leq p.
\]
Finally, Lemma~\ref{lem:excess-controls}, applied with \(A=M_F\) and
\(a=m\), gives
\[
s_r(F)
\leq rm+\sum_{\lambda\in\Spec(M_F)}(\lambda-m)_+
=rm+\tr((T_F)_+)
\leq rm+p_1
\leq rm+p,
\]
as required.
\end{proof}
Before proceeding to the main analysis, we record a global upper bound that
naturally applies when the condition $m < q$ is met.
\begin{lemma}\label{lem:low-core-bounds}
	 If $m<q,$ then  $s_r(F)\le D_r(F)\le p+rm$ for every \(r\ge1\).
\end{lemma}

\begin{proof}
	For every $r\ge1$,
	\[
	D_r(F)=p+\sum_{v\in Q}\min\{d_F(v),r\}\le p+rm.
	\]
	It remains to prove $s_r(F)\le D_r(F)$.
	Put $n=|F|$. Since $\sum_v d_F(v)=qn$ and
	$\min\{a,r\}=a-(a-r)_+$, we have
	\[
	D_r(F)=qn-\sum_{v\in Q}(d_F(v)-r)_+.
	\]
	If $r\ge n$, then $s_r(F)=\tr M_F=qn=D_r(F)$. Hence assume $r<n$.
	Since $d_F(v)\le n$, $D_r(F)\ge qn-m(n-r).$
	If $m\le q-2$, every facet is heavy, so
	Lemma~\ref{lem:heavy-isolated} gives $M_F=qI_n$. Thus
	\[
	s_r(F)=qr\le qr+(q-m)(n-r)=qn-m(n-r)\le D_r(F).
	\]
	Now let $m=q-1$, and let $b$ and $e$ be the numbers of one-pendant
	and heavy facets, respectively, so that $n=b+e$. If $b=0$, then
	\[
	s_r(F)=qr=rm+r\le rm+n.
	\]
	If $b\ge1$, Lemmas~\ref{lem:star} and \ref{lem:heavy-isolated} give
	the eigenvalues $b+m$, $m+1$ with multiplicity $e$, and $m$ with
	multiplicity $b-1$. Hence
	\[
	s_r(F)
	=rm+b+\min\{e,r-1\}
	\le rm+n.
	\]
	In either case,
	\[
	s_r(F)\le rm+n=qn-m(n-r)\le D_r(F).
	\]
\end{proof}
	
We call a \(q\)-family \emph{low-core} if
	\(|Q(F)|<q\), and \emph{high-core} if \(|Q(F)|\ge q\).
We now determine the precise equality conditions for the regime $m < q$.	

	\begin{lemma}\label{lem:low-equality}
		Let \(F\subseteq\binom Vq\) be a finite \(q\)-family without isolated vertices, where \(q\ge2\). Assume its core \(Q=Q(F)\) has size \(m=|Q|<q\).
		\begin{enumerate}[label=\textup{(\roman*)}]
			\item If \(m\le q-2\), then \(s_2(F)=D_2(F)\) if and only if \(|F|\le2\).
			\item If \(m=q-1\), then \(s_2(F)=D_2(F)\) if and only if \(F=\{Q\cup\{z\}:z\in Z\}\) with \(|Z|\ge2\) and \(Q\cap Z=\varnothing\).
		\end{enumerate}
	\end{lemma}
\begin{proof}
	Put \(P=P(F)\) and \(p=|P|\). Then \(D_2(F)=p+2m\) by \eqref{eq:R2-pq}.
Suppose \(m\le q-2\). Lemma~\ref{lem:heavy-isolated} implies \(M_F=qI_{|F|}\), whence \(s_2(F) = q\min\{|F|,2\}\). If \(|F|\le 2\), the core is the common intersection of all facets, which forces \(p=|F|(q-m)\) and yields \(D_2(F)=|F|q=s_2(F)\). Conversely, if \(|F|\ge 3\), the lower bound \(p\ge|F|(q-m)\) ensures
	\[
	D_2(F) \ge |F|(q-m)+2m \ge 3q-m > 2q = s_2(F),
	\]
	ruling out equality.
Now suppose \(m=q-1\). Let \(b\) and \(e\) denote the number of one-pendant and heavy facets in \(F\), respectively. If \(e=0\), then \(F\) is a pure star over \(Q\) with \(b\ge 2\) to satisfy the core degree constraints, and Lemma~\ref{lem:star} immediately yields the equality \(s_2(F) = (b+q-1)+(q-1) = b+2q-2 = D_2(F)\).
	
	It remains to show strict inequality whenever \(e\ge 1\). Since each heavy facet contributes at least two pendant vertices, we have \(p\ge b+2e\). If \(b\ge 1\), Lemmas~\ref{lem:star} and \ref{lem:heavy-isolated} yield \(s_2(F)=b+2q-1\), whereas
	\[
	D_2(F) = p+2q-2 \ge b+2e+2q-2 \ge b+2q > s_2(F).
	\]
	If \(b=0\), Lemma~\ref{lem:heavy-isolated} gives \(M_F=qI_{|F|}\). The condition \(d_F(v)\ge 2\) for all \(v\in Q\) forces \(|F|=e\ge 2\), implying \(p\ge 2e\ge 4\). Consequently,
	\[
	s_2(F) = 2q < p+2q-2 = D_2(F),
	\]
	which completes the proof.
\end{proof}

Fix \(q\geq2\) and a finite set \(Q\) with \(m=|Q|\geq q\). Retain the
complete-simplex notation \(\calF_Q\), \(\calR_Q\), \(\bd_Q\), and
\(\scrH_Q\) introduced above. Let \(\calK\subseteq\calF_Q\), put
\(\calN=\calF_Q\setminus\calK\), and define
\[
\scrV_{\calK}
=
\{\mathbf x\in\scrH_Q:\supp(\mathbf x)\subseteq\calK\}.
\]
Thus \(\scrV_{\calK}\) is the subspace of the complete-simplex row
space supported on \(\calK\).
We next evaluate the exact dimension of $\scrV_{\calK}$ and record its structural implications for $\calK$.
\begin{lemma}\label{lem:support-rank}
The subspace \(\scrV_{\calK}\) has dimension
\[
\dim\scrV_{\calK}
=
\binom{m-1}{q-1}-\rank_{\mathbb R}\bd_{\calN}.
\]
Moreover, if \(\dim\scrV_{\calK}\geq2\), then every vertex of \(Q\)
belongs to at least two members of \(\calK\).
\end{lemma}

\begin{proof}
Put \(h=\binom{m-1}{q-1}\). Let \(P_{\calN}:\scrH_Q\to\mathbb R^{\calN}\) be coordinate
restriction. If \(\mathbf x=\bd_Q^{\mathsf T}\mathbf z\), then, for
every \(\sigma\in\calN\),
\[
(P_{\calN}\mathbf x)_\sigma
=(\bd_Q^{\mathsf T}\mathbf z)_\sigma
=\langle\bd_Qe_\sigma,\mathbf z\rangle
=(\bd_{\calN}^{\mathsf T}\mathbf z)_\sigma.
\]
Thus \(P_{\calN}\bd_Q^{\mathsf T}=\bd_{\calN}^{\mathsf T}\).
Since \(\scrH_Q=\im\bd_Q^{\mathsf T}\), it follows that
\(\im P_{\calN}=\im\bd_{\calN}^{\mathsf T}\). Moreover,
\(P_{\calN}\mathbf x=0\) exactly when
\(\supp(\mathbf x)\subseteq\calK\), so
\(\Ker P_{\calN}=\scrV_{\calK}\). Since
\(\dim\scrH_Q=h\) by Lemma~\ref{lem:complete-simplex}, rank--nullity
gives
\[
\dim\scrV_{\calK}
=h-\rank P_{\calN}
=h-\rank_{\mathbb R}\bd_{\calN}.
\]
Fix \(v\in Q\) and put $\calF_v=\{\sigma\in\calF_Q:v\in\sigma\}.$
Every member of \(\calF_v\) has the unique form
\(\rho\cup\{v\}\), where
\(\rho\in\binom{Q\setminus\{v\}}{q-1}\). In
\(\bd_{\calF_v}\), restrict to the rows indexed by these sets \(\rho\).
The resulting square matrix is a signed identity: the row indexed by
\(\rho\) has a nonzero entry in the column indexed by
\(\rho\cup\{v\}\) and zero entries in all other selected columns.
Therefore $\rank\bd_{\calF_v}=|\calF_v|=h.$

Let \(P_v:\scrH_Q\to\mathbb R^{\calF_v}\) be coordinate restriction.
Since \(\rank\bd_{\calF_v}=h=|\calF_v|\), the transpose
\(\bd_{\calF_v}^{\mathsf T}\) has full row rank. Hence, for every
\(\mathbf y\in\mathbb R^{\calF_v}\), there exists \(\mathbf z\) such
that \(\bd_{\calF_v}^{\mathsf T}\mathbf z=\mathbf y\). Setting
\(\mathbf x=\bd_Q^{\mathsf T}\mathbf z\in\scrH_Q\), we obtain
\(P_v\mathbf x=\mathbf y\). Thus \(P_v\) is surjective. Since
\(\dim\scrH_Q=\dim\mathbb R^{\calF_v}=h\), surjectivity also implies
injectivity.
Now let \(\mathbf x\in\scrV_{\calK}\). Since
\(\supp(\mathbf x)\subseteq\calK\), the restricted vector
\(P_v\mathbf x\) is supported on \(\calK\cap\calF_v\). Because \(P_v\)
is injective,
\[
\dim\scrV_{\calK}
=
\dim P_v(\scrV_{\calK})
\leq|\calK\cap\calF_v|.
\]
Consequently, if \(\dim\scrV_{\calK}\geq2\), then
\(|\calK\cap\calF_v|\geq2\). Since \(v\) was arbitrary, every vertex
of \(Q\) belongs to at least two members of \(\calK\).
\end{proof}
	
Equipped with the preceding dimension constraints and core properties, we now proceed to the proof of Theorem~\ref{thm:main-inequality}.
\begin{proof}[\textup{\textbf{Proof of Theorem~\ref{thm:main-inequality}}}]
The case \(q=1\), including the empty family, was proved in Section~\ref{sec:notation} and gives case~\textup{(i)}. Assume \(q\geq2\), delete isolated vertices, and put \(P=P(F)\), \(Q=Q(F)\), \(p=|P|\), and \(m=|Q|\). Then \(D_2(F)=p+2m\) by \eqref{eq:R2-pq}. If \(m<q\), Lemmas~\ref{lem:low-core-bounds} and \ref{lem:low-equality} prove the inequality and show that equality holds exactly in cases~\textup{(ii)} and~\textup{(iii)}.

Assume \(m\geq q\), and put \(\widehat F=F\cup\binom Qq\), \(t=|\calC|\), and \(h=\binom{m-1}{q-1}\). For \(\rho\in\calC\), define
\[
Z_\rho=\{z\in P:\rho\cup\{z\}\in F\},\qquad b_\rho=|Z_\rho|,\qquad
\boldsymbol{\ell}_\rho=b_\rho^{-1/2}\sum_{z\in Z_\rho}\varepsilon(\rho,\rho\cup\{z\})e_{\rho\cup\{z\}},
\]
and put \(p_1=\sum_{\rho\in\calC}b_\rho\) and \(\mathbf u_\rho=(\mathbf y_\rho,\sqrt{b_\rho}\,\boldsymbol{\ell}_\rho)\).  Recall that
\[
\mathbf y_\rho=\bd_Q^{\mathsf T}e_\rho,
\qquad
(\mathbf y_\rho)_\sigma=\varepsilon(\rho,\sigma)
\quad(\sigma\in\calF_Q).
\]
When \(\mathbf y_\rho\) is viewed as a vector indexed by \(\widehat F\),
we set its coordinates outside \(\calF_Q\) equal to zero. Every facet
in \(\widehat F\setminus F\) belongs to \(\binom Qq\), so it contains
only core vertices. Thus completion adds no new one-pendant facets, and
\(\calC,Z_\rho,b_\rho,p_1\), and
\(\boldsymbol{\ell}_\rho\) remain unchanged.
 Since \(p_1\) counts the pendant vertices lying in one-pendant facets, \(p_1\leq p\).
 Lemmas~\ref{lem:completion} and~\ref{lem:core-excess} give
\begin{equation}\label{eq:main-high-chain}
s_2(F)\leq s_2(\widehat F)\leq2m+\tr((T_{\widehat F})_+)
\leq2m+p_1\leq2m+p=D_2(F).
\end{equation}
This proves the inequality.

Suppose that \(s_2(F)=D_2(F)\). Equality then holds throughout \eqref{eq:main-high-chain}. Thus \(p_1=p\), so there are no heavy facets. Equality in the positive-excess bound and Lemma~\ref{lem:positive-excess-rigidity} imply that the vectors \(\mathbf y_\rho\), \(\rho\in\calC\), are pairwise orthogonal and that \(M_{\widehat F}\mathbf u_\rho=(m+b_\rho)\mathbf u_\rho\). By Lemma~\ref{lem:complete-simplex}, pairwise orthogonality is equivalent to \(|\rho\cap\gamma|\leq q-3\) for distinct \(\rho,\gamma\in\calC\). If \(t\geq3\), writing \(b_1\geq\cdots\geq b_t>0\), we would have
\[
s_2(\widehat F)=2m+b_1+b_2<2m+\sum_{i=1}^t b_i=2m+p_1,
\]
contrary to \eqref{eq:main-high-chain}. Hence \(t\leq2\).
The equality statement of
Lemma~\ref{lem:positive-excess-rigidity}, together with the orthogonal
decomposition of \(M_{\widehat F}\) established in the proof of
Lemma~\ref{lem:core-excess}, gives
\[
\scrE_{>m}
=
\Span\{\mathbf u_\rho:\rho\in\calC\},
\qquad
\scrE_{=m}
=
\scrH_Q\cap
\Span\{\mathbf y_\rho:\rho\in\calC\}^{\perp}.
\]
Since the \(t\) vectors \(\mathbf y_\rho\) are nonzero and pairwise
orthogonal in the \(h\)-dimensional space \(\scrH_Q\), their orthogonal
complement in \(\scrH_Q\), namely \(\scrE_{=m}\), has dimension \(h-t\).
The same orthogonal decomposition shows that all remaining eigenvalues
of \(M_{\widehat F}\) are \(q-1\) or \(0\).  If \(h=1\), pairwise orthogonality forces \(t\leq1\); hence the largest eigenvalue is at most \(m+p\), and every subsequent eigenvalue is at most \(q-1<m\). This would give
\[
s_2(\widehat F)\leq(m+p)+(q-1)<2m+p,
\]
again contradicting \eqref{eq:main-high-chain}. Therefore \(h\geq2\), equivalently \(m\geq q+1\).

Since \(s_2(F)=s_2(\widehat F)\), Lemma~\ref{lem:compression-equality} gives a maximizing two-dimensional subspace \(\scrU\) of \(M_{\widehat F}\) supported on \(F\). Lemma~\ref{lem:kyfan} and the displayed eigenspaces show that, if \(t=0\), then \(\scrU\subseteq\scrH_Q\); if \(t=1\), with \(\calC=\{\rho\}\), then \(\scrU=\Span\{\mathbf u_\rho,\mathbf w\}\) for some nonzero \(\mathbf w\in\scrH_Q\cap\mathbf y_\rho^\perp\); and if \(t=2\), with \(\calC=\{\rho,\gamma\}\), then \(\scrU=\Span\{\mathbf u_\rho,\mathbf u_\gamma\}\). Because \(\scrU\) is supported on \(F\), and every \(\boldsymbol{\ell}_\rho\) is already supported on \(F\), these three cases give
\[
\dim\scrV_{\calK}\geq2,\qquad
\supp(\mathbf y_\rho)\subseteq\calK\quad(\rho\in\calC).
\]
Lemmas~\ref{lem:support-rank} and \ref{lem:complete-simplex} give
\[
\rank_{\R}\bd_{\calN}
\leq\binom{m-1}{q-1}-2,
\qquad
\rho\cup\{x\}\in\calK
\quad(\rho\in\calC,\ x\in Q\setminus\rho).
\]
We have also proved that \(m\geq q+1\), \(t=|\calC|\leq2\), and
$|\rho\cap\gamma|\leq q-3\ (\rho,\gamma\in\calC,\ \rho\neq\gamma).$
Finally, the absence of heavy facets means that every facet contains at
most one pendant vertex. Thus all the conditions in
case~\textup{(iv)} are satisfied.

Conversely, cases~\textup{(i)}--\textup{(iii)} were settled above;
assume now that case~\textup{(iv)} holds.  Condition~\textup{(a)} and  Lemma~\ref{lem:support-rank}   give \(\dim\scrV_{\calK}\geq2\), while condition~\textup{(c)} and Lemma \ref{lem:complete-simplex} give \(\mathbf y_\rho\in\scrV_{\calK}\) for every \(\rho\in\calC\). Condition~\textup{(b)} and
Lemma~\ref{lem:complete-simplex} show that these vectors are pairwise
orthogonal. There are no heavy facets, so \(p_1=p\); moreover, \(t\leq2\) and \(h\geq2\). Lemma~\ref{lem:positive-excess-rigidity}, together with the orthogonal
decomposition of \(M_{\widehat F}\) established in the proof of
Lemma~\ref{lem:core-excess}, shows that the two largest eigenvalues of
\(M_{\widehat F}\) are \(m,m\) when \(t=0\),
\(m+b_\rho,m\) when \(t=1\), and
\(m+b_\rho,m+b_\gamma\) when \(t=2\).  Hence
\[
s_2(\widehat F)=2m+\sum_{\rho\in\calC}b_\rho=2m+p=D_2(F).
\]
It remains to recover this value after deleting the missing core facets. If \(t=0\), choose any two-dimensional \(\scrU\subseteq\scrV_{\calK}\). If \(t=1\), choose a nonzero \(\mathbf w\in\scrV_{\calK}\cap\mathbf y_\rho^\perp\) and set \(\scrU=\Span\{\mathbf u_\rho,\mathbf w\}\). If \(t=2\), set \(\scrU=\Span\{\mathbf u_\rho,\mathbf u_\gamma\}\). In every case, \(\scrU\) is a maximizing two-dimensional subspace of \(M_{\widehat F}\) supported on \(F\). Lemma~\ref{lem:compression-equality} therefore gives
\[
s_2(F)=s_2(\widehat F)=D_2(F).
\]
This completes the equality classification.
\end{proof}
	
	

	\appendix
	\section{Exact Verification of the Seed Matrices}
	\label{app:seed-verification}
	
	This appendix records the integer Gram matrices underlying
	Lemma~\ref{lem:seed-polynomials}.  It is included to make the two seed
	calculations directly reproducible.  The vertex order is
	\(0<1<\cdots<6\), the facets are ordered exactly as below, and every facet
	is given its increasing orientation.  Reversing facet orientations only
	conjugates the corresponding Gram matrix by a diagonal sign matrix, so it
	does not affect the characteristic polynomial.
	
	For \(\calF_{15}\), use the facet order
	\[
	023,025,026,035,036,123,125,126,135,136,
	234,235,236,256,345.
	\]
We have
\begingroup
\normalsize
\setlength{\arraycolsep}{1.7pt}
\[
M_{\calF_{15}}=
\left(\begin{array}{*{15}{c}}
3& 1& 1&-1&-1& 1& 0& 0& 0& 0& 1& 1& 1& 0& 0\\
		1& 3& 1& 1& 0& 0& 1& 0& 0& 0& 0&-1& 0& 1& 0\\
		1& 1& 3& 0& 1& 0& 0& 1& 0& 0& 0& 0&-1&-1& 0\\
		-1& 1& 0& 3& 1& 0& 0& 0& 1& 0& 0& 1& 0& 0&-1\\
		-1& 0& 1& 1& 3& 0& 0& 0& 0& 1& 0& 0& 1& 0& 0\\
		1& 0& 0& 0& 0& 3& 1& 1&-1&-1& 1& 1& 1& 0& 0\\
		0& 1& 0& 0& 0& 1& 3& 1& 1& 0& 0&-1& 0& 1& 0\\
		0& 0& 1& 0& 0& 1& 1& 3& 0& 1& 0& 0&-1&-1& 0\\
		0& 0& 0& 1& 0&-1& 1& 0& 3& 1& 0& 1& 0& 0&-1\\
		0& 0& 0& 0& 1&-1& 0& 1& 1& 3& 0& 0& 1& 0& 0\\
		1& 0& 0& 0& 0& 1& 0& 0& 0& 0& 3& 1& 1& 0& 1\\
		1&-1& 0& 1& 0& 1&-1& 0& 1& 0& 1& 3& 1&-1&-1\\
		1& 0&-1& 0& 1& 1& 0&-1& 0& 1& 1& 1& 3& 1& 0\\
		0& 1&-1& 0& 0& 0& 1&-1& 0& 0& 0&-1& 1& 3& 0\\
		0& 0& 0&-1& 0& 0& 0& 0&-1& 0& 1&-1& 0& 0& 3
\end{array}\right).
\]
\endgroup
	
	For \(\calF_{16}\), use the facet order
	\[
	012,013,015,026,123,124,125,126,
	134,135,136,145,146,156,236,256.
	\]
	Then
\begingroup
\normalsize
\setlength{\arraycolsep}{1.7pt}
\[
M_{\calF_{16}}=
\left(\begin{array}{*{16}{c}}
3& 1& 1&-1& 1& 1& 1& 1& 0& 0& 0& 0& 0& 0& 0& 0\\
		1& 3& 1& 0&-1& 0& 0& 0& 1& 1& 1& 0& 0& 0& 0& 0\\
		1& 1& 3& 0& 0& 0&-1& 0& 0&-1& 0&-1& 0& 1& 0& 0\\
		-1& 0& 0& 3& 0& 0& 0& 1& 0& 0& 0& 0& 0& 0&-1&-1\\
		1&-1& 0& 0& 3& 1& 1& 1&-1&-1&-1& 0& 0& 0& 1& 0\\
		1& 0& 0& 0& 1& 3& 1& 1& 1& 0& 0&-1&-1& 0& 0& 0\\
		1& 0&-1& 0& 1& 1& 3& 1& 0& 1& 0& 1& 0&-1& 0& 1\\
		1& 0& 0& 1& 1& 1& 1& 3& 0& 0& 1& 0& 1& 1&-1&-1\\
		0& 1& 0& 0&-1& 1& 0& 0& 3& 1& 1&-1&-1& 0& 0& 0\\
		0& 1&-1& 0&-1& 0& 1& 0& 1& 3& 1& 1& 0&-1& 0& 0\\
		0& 1& 0& 0&-1& 0& 0& 1& 1& 1& 3& 0& 1& 1& 1& 0\\
		0& 0&-1& 0& 0&-1& 1& 0&-1& 1& 0& 3& 1&-1& 0& 0\\
		0& 0& 0& 0& 0&-1& 0& 1&-1& 0& 1& 1& 3& 1& 0& 0\\
		0& 0& 1& 0& 0& 0&-1& 1& 0&-1& 1&-1& 1& 3& 0& 1\\
		0& 0& 0&-1& 1& 0& 0&-1& 0& 0& 1& 0& 0& 0& 3& 1\\
		0& 0& 0&-1& 0& 0& 1&-1& 0& 0& 0& 0& 0& 1& 1& 3
\end{array}\right).
\]
\endgroup

The characteristic polynomials are computed directly from the two
displayed integer matrices. Applying Bareiss fraction-free elimination \cite{Bareiss1968}
to \(xI-M_{\calF}\) over \(\mathbb Z[x]\) gives
\[
\begin{aligned}
\det(xI-M_{\calF_{16}})
&=x^2(x-1)^5(x-3)(x-7)^3p_{16}(x),\\
\det(xI-M_{\calF_{15}})
&=x^4(x-2)(x-4)^2(x-6)^2(x-7)p_{15}(x).
\end{aligned}
\]
All divisions in the Bareiss elimination are exact by the
Desnanot--Jacobi identity, and the pivots are nonzero because the
leading principal minors of \(xI-M_{\calF}\) are monic polynomials.
Thus the computation uses exact polynomial arithmetic and no numerical
eigenvalue approximation.


\newpage
 

{\bf Declaration on the use of AI}

During the preparation of this manuscript, we used ChatGPT and DeepSeek  for language editing,
which included grammar verification, stylistic polishing, and the improvement of expository clarity.
	
\begin{thebibliography}{99}

\bibitem{Abebe}
R. Abebe, A conjectural Brouwer inequality for higher-dimensional Laplacian spectra, arXiv
(2019),  arXiv:1907.07541.

\bibitem{Abebe12}
R. Abebe and J. Pfeffer, On the partial sum of the Laplacian eigenvalues of abstract
simplicial complexes. Unpublished report, https://www-users.cse.umn.edu/~reiner/REU/
AbebePfeffer2012.pdf.

	\bibitem{Bai2011}
	H. Bai,
	\newblock The Grone--Merris conjecture,
	\newblock \emph{Transactions of the American Mathematical Society}
	\textbf{363} (2011), no.~8, 4463--4474.
	

\bibitem{BapatLalPati2008}
	R. B. Bapat, A. K. Lal, and S. Pati,
	\newblock Laplacian spectrum of weakly quasi-threshold graphs,
	\newblock \emph{Graphs and Combinatorics} \textbf{24} (2008),
	273--290.

	\bibitem{Bareiss1968}
	E. H. Bareiss,
	\newblock Sylvester's identity and multistep integer-preserving Gaussian
	elimination,
	\newblock \emph{Mathematics of Computation} \textbf{22} (1968),
	no.~103, 565--578.

	\bibitem{Bhatia}
	R. Bhatia,
	\newblock \emph{Matrix Analysis},
	\newblock Graduate Texts in Mathematics, Vol.~169, Springer, New York,
	1997.

	\bibitem{Duval2005}
	A. M. Duval,
	\newblock A common recursion for Laplacians of matroids and shifted
	simplicial complexes,
	\newblock \emph{Documenta Mathematica} \textbf{10} (2005), 583--618.
	

	\bibitem{DuvalReiner2002}
	A. M. Duval and V. Reiner,
	\newblock Shifted simplicial complexes are Laplacian integral,
	\newblock \emph{Transactions of the American Mathematical Society}
	\textbf{354} (2002), no.~11, 4313--4344.
	

	\bibitem{Eckmann1944}
	B. Eckmann,
	\newblock Harmonische Funktionen und Randwertaufgaben in einem Komplex,
	\newblock \emph{Commentarii Mathematici Helvetici} \textbf{17}
	(1944), 240--255.
	
	\bibitem{Fan}
	K. Fan,
	\newblock On a theorem of Weyl concerning eigenvalues of linear
	transformations I,
	\newblock \emph{Proceedings of the National Academy of Sciences of the
	United States of America} \textbf{35} (1949), no.~11, 652--655.
	

	\bibitem{FanWuWang2025}
	Y. Fan, H.  Wu, and Y. Wang,
	\newblock The largest Laplacian eigenvalue and the balancedness of
	simplicial complexes,
	\newblock \emph{Journal of Algebraic Combinatorics} \textbf{61}
	(2025), Article~53.
	

	\bibitem{GroneMerris1994}
	R. Grone and R. Merris,
	\newblock The Laplacian spectrum of a graph. II,
	\newblock \emph{SIAM Journal on Discrete Mathematics} \textbf{7}
	(1994), no.~2, 221--229.
	

	\bibitem{HanLuWang2026}
	Y. Han, L. Lu, and J. Wang,
	\newblock On the sum of the two largest eigenvalues of the curl-curl
	operator on graphs,
	arXiv (2026),
arXiv:2606.26512v1.

	\bibitem{Katz2005}
	N. H. Katz,
	\newblock The Grone--Merris conjecture and a quadratic eigenvalue
	problem,
	arXiv (2005),
arXiv:math/0512647.

	\bibitem{Kirchhoff1847}
	G. Kirchhoff,
	\newblock Ueber die Aufl\"osung der Gleichungen, auf welche man bei der
	Untersuchung der linearen Vertheilung galvanischer Str\"ome gef\"uhrt
	wird,
	\newblock \emph{Annalen der Physik und Chemie} \textbf{148} (1847),
	no.~12, 497--508.

	\bibitem{Kirkland2009}
	S. Kirkland,
	\newblock Near threshold graphs,
	\newblock \emph{The Electronic Journal of Combinatorics} \textbf{16}
	(2009), no.~1, Research Paper~R42.

	\bibitem{Lew2025}
	A. Lew,
	\newblock Sums of Laplacian eigenvalues and sums of degrees,
 arXiv (2025),
arXiv:2508.04209v1.

	\bibitem{Stephen2007}
	T. Stephen,
	\newblock A majorization bound for the eigenvalues of some graph
	Laplacians,
	\newblock \emph{SIAM Journal on Discrete Mathematics} \textbf{21}
	(2007), no.~2, 303--312.
	

	\bibitem{ZhangFan2026}
	H. Zhang and Y.  Fan,
	\newblock Upper bounds for the largest Laplacian eigenvalue of simplicial
	complexes, arXiv  (2026),
arXiv:2606.21233v1.

\end{thebibliography}


	
\end{document}
